\subsection{Local full-to-LOO stability with an empirical degree floor}
\label{restore:loo}

The row bound used in the dense-range LOO argument need only hold on
the support of the tested row. This distinction removes a union over
all vertices. The threshold is recomputed after deletion.

\begin{lemma}[A local floored scalar comparison]\label{restore:loo-scalar}
Write $A=\left(\begin{smallmatrix}0&x^\top\\x&B\end{smallmatrix}\right)$,
$k=\sum_jx_j\le Ld$, and let $\Delta$ be the diagonal degree floor
recomputed from $B$. Suppose
\[
 cd\le\Delta_j,\qquad \sum_j\Delta_j\le Cnd,
 \qquad S_B=\Delta^{-1/2}B\Delta^{-1/2}
   =U_B\Lambda_BU_B^\top+N,
\]
where $NU_B=0$, $\norm N_{\op}\le C/\sqrt d$, and the $K$ retained
eigenvalues have a fixed gap from zero. Let the full floor be
$G=\diag(g,\Delta+E)$, with
\[
 cd\le g\le Cd,\qquad E=E_\mathrm{loc}+E_\mathrm{glob},\qquad
 0\le E_\mathrm{loc}\le\diag(x),\quad
 \norm{E_\mathrm{glob}}_{\op}\le Cd/n.
\]
Here $d\to\infty$ and $d\le Cn$. Suppose only on hit coordinates that
\begin{equation}\label{restore:loo-hit}
 \norm{\Delta_j^{-1/2}(U_B)_{j\cdot}}_2
       \le R/\sqrt{nd}\quad\text{whenever }x_j=1.
\end{equation}
Let $U,\Lambda$ be the full informative eigensystem of
$G^{-1/2}AG^{-1/2}$, and put
\[
 H_G=G^{-1/2}UU^\top G^{1/2},\qquad
 H_B=\Delta^{-1/2}U_BU_B^\top\Delta^{1/2}.
\]
For every $\norm w_\infty\le W$, uniformly over the stated inputs,
\begin{equation}\label{restore:loo-main}
 (AH_Gw)_i=x^\top H_Bw_{-i}+O_{L,R,W}(1).
\end{equation}
No bound is assumed on unhit rows of $U_B$ or on $\max_j\Delta_j$.
\end{lemma}

\begin{proof}
For $k=0$ both row actions vanish. For $k>0$, put
$F=\Delta^{-1/2}E\Delta^{-1/2}$, $h=\Delta^{-1/2}x$,
$Q=I-U_BU_B^\top$. Then $\norm F\le C/d$ and
$\norm h^2\le Ck/d$. The full normalized matrix differs from
$\diag(0,S_B)$ by $O(d^{-1/2})$: its off-diagonal row has norm
at most $C\sqrt{k/(gd)}$, and the lower renormalization has norm
$O(d^{-1})$. Thus the full signal remains separated.

Write a full generalized eigenvector as
\[
 G^{-1/2}u_r=\binom{t_r}{\Delta^{-1/2}(U_By_r+z_r)},
 \qquad U_B^\top z_r=0.
\]
The exact top and complementary equations are
\begin{align*}
 h^\top(U_By_r+z_r)&=\mu_r g t_r,\\
 (\mu_rI-N+\mu_rQFQ)z_r
   &=t_rQh-\mu_rQFU_By_r.
\end{align*}
The inverse $T_r$ of the complementary operator is bounded. Therefore
\[
 (\mu_rg-(Qh)^\top T_rQh)t_r
  =(U_B^\top h)^\top y_r
       -\mu_r(Qh)^\top T_rQFU_By_r.
\]
The denominator has magnitude at least $cd-Ck/d\ge c'd$.
It is the floored degree $g$, not the raw degree $k$, in this formula.
The matrix $Y=(y_1,\ldots,y_K)$ and its inverse have bounded norm,
by the signal-space angle and $(I+F)^{1/2}=I+O(d^{-1})$.

The local row premise and the global threshold perturbation give
\[
 \norm{U_B^\top h}\le C_Rk/\sqrt{nd},\qquad
 \norm{FU_B}_F\le C_R\sqrt{k}/(d\sqrt n)+C/n.
\]
The $C/n$ term is retained; it need not be supported on the neighbors.
Substitution yields, for $Z_\perp=(z_1,\ldots,z_K)$,
\begin{equation}\label{restore:loo-refined}
 \norm t\le C_R/\sqrt{nd},\qquad
 \norm{Z_\perp}\le C_R\{\sqrt{k}/(d\sqrt n)+1/n\}.
\end{equation}
Let $V=\Delta^{-1/2}U_B$, $v_r=\Delta^{-1/2}z_r$ and
$f_r=Vy_r+v_r$. Exact generalized orthonormality gives
\[
 Y^\top Y+Z_\perp^\top Z_\perp+gtt^\top+f^\top Ef=I.
\]
On the local part of $E$, use the hit-row bound; on its global part
use $\norm{E_\mathrm{glob}}\le Cd/n$ and
$\norm f\le C/\sqrt d$. Consequently
\begin{equation}\label{restore:loo-gram}
 \norm{YY^\top-I}\le C_R\{k/(nd)+1/n\}.
\end{equation}

For $w_V=V^\top\Delta w_{-i}$, the exact row expansion is
\[
 (AH_Gw)_i=
 \sum_r(x^\top Vy_r+x^\top v_r)
 \{g t_rw_i+y_r^\top w_V+v_r^\top\Delta w_{-i}
                                      +f_r^\top Ew_{-i}\}.
\]
Only total degree is needed to bound global coefficient vectors. In
particular,
\begin{align*}
 \norm{x^\top V}&\le C_Rk/\sqrt{nd},&
 \norm{w_V}&\le C_W\sqrt{nd},\\
 \norm{x^\top v}&\le C_R\{k/(d^{3/2}\sqrt n)+\sqrt{k}/(n\sqrt d)\},&
 \norm{v^\top\Delta w_{-i}}&\le C_{R,W}\{\sqrt{k/d}+\sqrt{d/n}\},\\
 \norm{gt}&\le C_R\sqrt{d/n},&
 \norm{f^\top Ew_{-i}}&\le C_{R,W}\{k/\sqrt{nd}+\sqrt{d/n}\}.
\end{align*}
For example the global part of the last bound follows from
$|f_r^\top E_\mathrm{glob}w_{-i}|
\le (Cd/n)W\sqrt n\norm{f_r}\le C_W\sqrt{d/n}$.
The leading term is $x^\top VYY^\top w_V$; its error relative to
$x^\top VV^\top\Delta w_{-i}$ is
$O(k^2/(nd)+k/n)$. Every other product is bounded by a fixed
multiple of
\[
 k/d+\sqrt{k/n}+k^2/(nd)+k/n+k\sqrt{k}/(d\sqrt n).
\]
These quantities are bounded for $k\le Ld$ and $d\le Cn$.
This proves the scalar comparison.
\end{proof}

For the empirical floor $\tau=a\bar D$, the displayed decomposition of
$E$ is automatic on the two-sided total-degree event. Indeed, if
$\bar D_B$ is the mean degree of $B$, then
\[
 \tau_A-\tau_B=a(2k-\bar D_B)/n.
\]
First change each deleted degree by $x_j$ while keeping $\tau_B$;
this yields $E_\mathrm{loc}\in[0,\diag(x)]$. Then change the threshold;
this yields $\norm{E_\mathrm{glob}}\le C_Ld/n$.

\begin{corollary}[Exact-degree reconstruction and simple hollowing]
\label{restore:loo-corrected}
Put $R_0=AH_G$ and $\mathbf d=A\one$, and define
\begin{equation}\label{restore:reconstruction}
 R_1=R_0+(\mathbf d-R_0\one)\one^\top/n,\qquad
 R^h=R_1+\frac{\diag(R_1)}{n-1}\one^\top
                -\frac n{n-1}\diag\{\diag(R_1)\}.
\end{equation}
In this display $\diag(R_1)$ in the numerator denotes its diagonal
column vector; the final $\diag\{\cdot\}$ forms a diagonal matrix.
Then $R^h\one=\mathbf d$ and $\diag(R^h)=0$ exactly.
For $\bar w_{-i}=(n-1)^{-1}\sum_{j\ne i}w_j$, the same local premises
give
\begin{equation}\label{restore:loo-corrected-action}
 (R^h w)_i
 =x^\top\{H_B(w_{-i}-\bar w_{-i}\one)+\bar w_{-i}\one\}
       +O_{L,R,W}(1).
\end{equation}
\end{corollary}

\begin{proof}
For $\bar w=n^{-1}\sum_jw_j$, exact degree correction gives
$(R_1w)_i=(R_0(w-\bar w\one))_i+k\bar w$.
Apply the preceding lemma. On hit coordinates $|(H_B\one)_j|\le C_R$,
because $\norm{U_B^\top\Delta^{1/2}\one}\le C\sqrt{nd}$.
Replacing $\bar w$ by $\bar w_{-i}$ therefore costs $O(k/n)$.
Moreover \eqref{restore:loo-refined} gives
$\norm{u_i}^2=g\norm t^2\le C_R/n$, hence
$|(R_0)_{ii}|\le C_Rd/n$. Also $(R_0\one)_i=O_R(k+1)$.
Thus $|(R_1)_{ii}|\le C_Rd/n$ and the hollow change to a bounded
contrast is at most $2W|(R_1)_{ii}|=O(d/n)$. The two exact identities
follow by summing the defined rows and examining their diagonal.
\end{proof}

The independent coefficient vector in
\eqref{restore:loo-corrected-action} has perturbation
$(H_B-I)(w_{-i}-\bar w_{-i}\one)$. For block-constant contrasts it
therefore has the same tilted cancellation and quadratic bound already
proved in \eqref{eq:floor-tilted-cancellation}--\eqref{eq:floor-tilt-moments}.
For clarity, the hit-row exception is paid directly. On the
training event, $\Delta_j\ge c d$, so
\[
 \mathcal J_{\rm bad}
 :=\{j:\norm{\Delta_j^{-1/2}(U_B)_{j\cdot}}_2
                          >R/\sqrt{nd}\}
 \subseteq
 \{j:\sqrt n\norm{(U_B)_{j\cdot}}_2>R\sqrt c\}.
\]
\Cref{restore:tail-energy-lemma} bounds the expected cardinality of
this set on the training event by $Cn e^{-c'd\log R}$.
Conditioning only on the deleted graph and using $p_{ij}\le Cd/n$
therefore bounds the probability that the tested row hits it by
$Cd e^{-c'd\log R}$. Choosing a sufficiently large fixed $R$
gives any prescribed fixed exponential rate $H$.
The upper-degree exception $k>Ld$ has arbitrarily large exponential
rate after increasing the fixed $L$. Neither bound uses a union over
all possible tested targets.

Uniformity in $w$ is an algebraic stability statement. It does not make
full-graph fitted coefficients independent of $x$; the following fitted
parameter lemma replaces them by deterministic true contrasts before
the independent sharp-tail argument is used.
